
\documentclass[11pt]{article}
\usepackage{geometry}
\usepackage{hyperref}
\usepackage{longtable}
\usepackage{graphicx}
\usepackage{listings}
\usepackage{float}
\geometry{margin=1in}

\title{Standard Operating Procedure (SOP)\\Wastewater Pathogen Sequencing Analysis Pipeline}
\author{Dieter Best}
\date{\today}

\begin{document}
\maketitle

\section{Purpose}
This document describes the Standard Operating Procedure (SOP) for running the wastewater pathogen sequencing pipeline. The pipeline processes Illumina paired-end sequencing data from wastewater samples to:

\begin{itemize}
\item Perform read quality control
\item Remove contaminants and sequencing artifacts
\item Classify microbial reads
\item Align reads to a pathogen reference genome
\item Perform primer trimming and variant analysis
\item Estimate lineage composition
\item Generate coverage and quality metrics
\item Produce a consolidated MultiQC report
\end{itemize}

The pipeline is implemented in Workflow Description Language (WDL) and executed using workflow engines such as Cromwell.

\section{Scope}

This pipeline is designed for analysis of:

\begin{itemize}
\item Wastewater sequencing samples
\item Illumina paired-end FASTQ data
\item Amplicon-based pathogen sequencing
\end{itemize}

Typical targets include:

\begin{itemize}
\item SARS-CoV-2
\item Influenza
\item RSV
\item Other viral pathogens
\end{itemize}

\section{Pipeline Overview}

The workflow performs the following major stages:

\begin{enumerate}
\item Raw read quality control
\item Read trimming and contaminant removal
\item Optional taxonomic classification
\item Alignment to pathogen reference genome
\item Alignment quality assessment
\item Primer trimming
\item Coverage calculation
\item Variant and lineage analysis
\item Aggregated reporting
\end{enumerate}

Each sample is processed independently using WDL scatter parallelization.

\section{Input Data}

\subsection{Sequencing Reads}

\begin{longtable}{p{4cm} p{10cm}}
\textbf{Input} & \textbf{Description} \\
reads1 & Forward FASTQ files \\
reads2 & Reverse FASTQ files \\
samplenames & Sample identifiers \\
\end{longtable}

\subsection{Reference Data}

\begin{longtable}{p{4cm} p{10cm}}
\textbf{Input} & \textbf{Description} \\
reference & Pathogen reference genome \\
primers & Primer sequences \\
primers\_bed & Primer coordinates \\
\end{longtable}

\subsection{Adapter and Contamination Databases}

\begin{longtable}{p{4cm} p{10cm}}
\textbf{File} & \textbf{Purpose} \\
adapters & Adapter sequences \\
phiX & PhiX contamination removal \\
polyA & Poly-A removal \\
\end{longtable}

\subsection{Taxonomic Classification (Optional)}

\begin{longtable}{p{4cm} p{10cm}}
\textbf{Input} & \textbf{Description} \\
run\_centrifuge & Enable classification \\
indexFiles & Centrifuge database \\
\end{longtable}

\subsection{Variant and Lineage Analysis}

\begin{longtable}{p{4cm} p{10cm}}
\textbf{Input} & \textbf{Description} \\
pathogen & Pathogen name \\
depth\_cut\_off & Minimum depth for lineage analysis \\
min\_base\_quality & Minimum base quality \\
\end{longtable}

\section{Pipeline Steps}

\subsection{Read Quality Control}

\textbf{Tool:} FastQC

Purpose:

\begin{itemize}
\item Evaluate base quality scores
\item Assess GC content
\item Detect adapter contamination
\item Identify overrepresented sequences
\end{itemize}

Outputs include HTML reports and QC metrics.

\subsection{Read Trimming and Filtering}

\textbf{Tool:} fastp

Purpose:

\begin{itemize}
\item Remove adapter sequences
\item Trim low-quality bases
\item Filter short reads
\end{itemize}

Outputs:

\begin{itemize}
\item Cleaned FASTQ files
\item JSON QC report
\item HTML QC report
\end{itemize}

\subsection{Contamination Removal}

\textbf{Tool:} BBDuk

Purpose:

Removes sequencing contaminants including:

\begin{itemize}
\item Adapter sequences
\item PhiX control reads
\item Poly-A artifacts
\item Primer sequences
\end{itemize}

Outputs cleaned FASTQ files and contamination statistics.

\subsection{Taxonomic Classification}

\textbf{Tool:} Centrifuge

This optional step classifies reads taxonomically to identify organisms present in wastewater samples.

Outputs include classification tables and summary statistics.

\subsection{Alignment to Pathogen Reference}

\textbf{Tool:} Minimap2

Reads are aligned to the pathogen reference genome to generate BAM alignment files.

\subsection{Alignment Quality Assessment}

\textbf{Tool:} Qualimap

Calculates:

\begin{itemize}
\item coverage statistics
\item insert size distribution
\item mapping quality
\item GC bias
\end{itemize}

\subsection{BAM Metrics Collection}

\textbf{Tools:} Picard and Samtools

Generates alignment summary metrics and quality-by-cycle metrics.

\subsection{Whole Genome Metrics}

\textbf{Tool:} GATK

Collects genome-wide coverage statistics.

\subsection{Coverage Calculation}

\textbf{Tool:} Mosdepth

Computes coverage depth across the genome and across targeted primer regions.

Outputs:

\begin{itemize}
\item Per-base coverage
\item Coverage summary
\item Coverage distribution
\end{itemize}

\subsection{Primer Trimming}

\textbf{Tool:} iVar

Removes primer sequences from aligned reads and prepares alignments for variant analysis.

\subsection{Lineage and Variant Deconvolution}

\textbf{Tool:} Freyja

Estimates lineage composition in mixed wastewater samples and identifies pathogen variants.

Outputs include lineage abundance estimates and variant calls.

\section{MultiQC Reporting}

All pipeline metrics are aggregated into a single report using MultiQC.

The report includes:

\begin{itemize}
\item FastQC metrics
\item fastp trimming statistics
\item alignment metrics
\item coverage statistics
\item lineage composition
\end{itemize}

\section{Pipeline Flow Diagram}

\begin{verbatim}
Raw FASTQ
   |
   v
FastQC
   |
   v
fastp trimming
   |
   v
BBDuk contaminant removal
   |
   +---- Centrifuge classification (optional)
   |
   v
Minimap2 alignment
   |
   v
Qualimap QC
   |
   v
Picard / WGS metrics
   |
   v
Mosdepth coverage
   |
   v
iVar primer trimming
   |
   v
Freyja lineage analysis
   |
   v
MultiQC report
\end{verbatim}

\section{Parallel Execution}

Samples are processed in parallel using WDL scatter blocks, allowing efficient processing of large sample batches.

\section{Containerization}

All tools are executed within Docker containers defined in the pipeline configuration. This ensures reproducibility and consistent software environments.

\section{Example Execution}

Example Cromwell command:

\begin{lstlisting}
java -jar cromwell.jar run wf_wastewater_pathogen.wdl \
  --inputs inputs.json
\end{lstlisting}

\section{Pipeline Outputs}

Major outputs include:

\begin{itemize}
\item FastQC reports
\item trimmed FASTQ files
\item BAM alignment files
\item coverage metrics
\item lineage estimates
\item MultiQC report
\end{itemize}

\end{document}
